\SourceLeaf{402}%
\noindent Diess substituirt, erhalten wir
\[
\begin{aligned}
&2(Pt't'+2Qt'\theta'+R\theta'\theta')\BareRoot{(PR-QQ)}.(df-d\chi)\\
&\quad=t'^3(-QP'+2PQ'-PP'')du\\
&\qquad+t't'\theta'(PR'-3QP''-\frac{2QQ}{P}P'+4QQ')du\\
&\qquad+t'\theta'\theta'(2QR'-\frac{2QQ}{P}P''-\frac{RQ}{P}P'+2RQ'-RP'')du\\
&\qquad+\theta'^3\left(RR'-\frac{QR}{P}P''\right)du.
\end{aligned}
\]
[Die rechte Seite ist aber identisch]
\[
\begin{aligned}
=(Pt't'+2Qt'\theta'+R\theta'\theta')\{&\left(-\frac QP P'+2Q'-P''\right)t'du\\
&+\left(R'-\frac QP P''\right)\theta'du\}.
\end{aligned}
\]

Wir haben also:
\begin{equation}\tag*{[22]}
\begin{aligned}
&2P\BareRoot{(PR-QQ)}(df-d\chi)\\
&\quad=\{(-QP'+2PQ'-PP'')t'+(PR'-QP'')\theta'\}du\\
&\quad=-QdP+(2PQ't'+PR'\theta')du-PP''t'du,
\end{aligned}
\end{equation}
[wofür man auch schreiben kann:]
\begin{equation}\tag*{[23]}
\begin{aligned}
&2\BareRoot{(PR-QQ)}(df-d\chi)\\
&\quad=\frac{-QP'+2PQ'-PP''}{P}dt-\frac{-RR'+2RQ''-QR''}{R}d\theta\\
&\qquad-Q\left(\frac{P''}{P}-2\frac{Q'}{Q}+\frac{R''}{R}\right)d\theta.
\end{aligned}
\end{equation}
\Needspace{8\baselineskip}
\begin{center}\bfseries Bemerkungen.\end{center}
\noindent\hspace*{1.2em}In \S\ 13 der \textit{Disquisitiones generales circa superficies curvas} stellt \textsc{Gauss} die Veröffentlichung weiterer Untersuchungen über die »absoluten« Eigenschaften krummer Flächen in Aussicht, die wie das Krümmungsmass und die geodätischen Linien gegenüber Biegungen der Flächen invariant sind. Als Vorarbeit hierzu ist die vorstehende Notiz aufzufassen, die höchst wahrscheinlich aus der Zeit zwischen 1822 und 1825 stammt. Sie ist auf einzelne Zettel geschrieben, augenscheinlich in grosser Eile. Daher kommt es wohl, dass \textsc{Gauss} die »Bestimmenden der Fläche« bald $t,u$, bald $t,\theta$ nennt, während hier durchgängig $t,\theta$ gesetzt worden ist; ebenso, dass er die Buchstaben $\sigma$ und $\psi$, die in den Gleichungen 3 und 9 vorkommen, später auch für die Grössen verwendet, die hier mit $\tau$ (Gleichung 6) und $\chi$ (Gleichung 17) bezeichnet worden sind. Auch sonst
